\documentclass[11pt,a4paper]{article}
\usepackage[T1]{fontenc}
\usepackage[utf8]{inputenc}
\usepackage{mathptmx}
\AtBeginDocument{\let\hbar\relax
  \DeclareMathSymbol{\hbar}{\mathord}{AMSb}{"7E}}
\usepackage{amsmath,amssymb,amsthm}
\usepackage{geometry}
\usepackage{microtype}
\usepackage{hyperref}
\usepackage{booktabs}
\usepackage{enumitem}
\usepackage{natbib}
\usepackage{titlesec}

\geometry{top=1.1in,bottom=1.1in,left=1.15in,right=1.15in}

\titleformat{\section}{\normalfont\bfseries\large}{\thesection}{1em}{}
\titleformat{\subsection}{\normalfont\bfseries\normalsize}{\thesubsection}{1em}{}
\titleformat{\subsubsection}{\normalfont\bfseries\itshape\normalsize}
  {\thesubsubsection}{1em}{}
\titlespacing*{\section}{0pt}{14pt}{4pt}
\titlespacing*{\subsection}{0pt}{10pt}{3pt}

\theoremstyle{plain}
\newtheorem{theorem}{Theorem}[section]
\newtheorem{proposition}[theorem]{Proposition}
\newtheorem{corollary}[theorem]{Corollary}
\theoremstyle{definition}
\newtheorem{definition}[theorem]{Definition}
\newtheorem{conjecture}[theorem]{Conjecture}
\newtheorem{postulate}[theorem]{Postulate}
\theoremstyle{remark}
\newtheorem{remark}[theorem]{Remark}

\hypersetup{colorlinks=true,linkcolor=black,citecolor=black,urlcolor=black,
  pdftitle={The Engine of Becoming (v2)},
  pdfauthor={Joshua F. Sandeman}}

\newcommand{\Pact}{\mathcal{P}_{\mathrm{act}}}
\newcommand{\Tact}{T^{\mu\nu}_{\mathrm{RA}}}
\newcommand{\ARA}{A_{\mathrm{RA}}}
\newcommand{\Vvalid}{\mathcal{V}_{\mathrm{valid}}}
\newcommand{\ket}[1]{|#1\rangle}
\newcommand{\muC}{\mu_{\mathrm{QM}}}   % quantum measure

% ─────────────────────────────────────────────────────────────────────────────
\begin{document}

% ── Title block ───────────────────────────────────────────────────────────────
\begin{center}
  {\LARGE\bfseries The Engine of Becoming:\par}
  \vspace{6pt}
  {\large\bfseries A Covariant Stochastic Graph Rewriting Algorithm as the\\
   Unified Foundation of Quantum Mechanics and General Relativity\par}
  \vspace{14pt}
  {\large Joshua F.\ Sandeman$^{1*}$\par}
  \vspace{4pt}
  {\normalsize $^{1}$Independent Researcher, Salem, OR, USA.\par}
  {\normalsize $^{*}$Corresponding author: sansuikyo@gmail.com\par}
  \vspace{8pt}
  {\small\textit{Preprint:}~\href{https://doi.org/10.5281/zenodo.19198120}
    {doi.org/10.5281/zenodo.19198120}\par}
  \vspace{14pt}
\end{center}

% ── Abstract ─────────────────────────────────────────────────────────────────
\begin{abstract}
The incompatibility of quantum mechanics and general relativity arises from a
shared foundational error: both frameworks treat time as a background parameter
rather than a physical process. Relational Actualism (RA) dissolves this tension
by replacing both the Schr\"{o}dinger equation and the Einstein field equations
with a single \textbf{Covariant Stochastic Graph Rewriting System}---the Engine
of Becoming. The algorithm operates on the growing causal Directed Acyclic Graph
(DAG) of actualization events at three levels: (1) a graph-theoretic primitive
layer describing candidate extensions of the causal boundary; (2) a quantum
measure layer assigning complex amplitudes to causal histories in a
foliation-independent way, using Sorkin's quantum measure formalism; and (3) an
emergent layer in which the Feynman path integral and the Einstein field equations
are the macroscopic descriptions: the quantum measure is exact at Level~2 and
the Feynman integral is its large-$\mu$ approximation; the Einstein field
equations are the unique consistent Level~3 description, inevitable by Lovelock's
theorem~\citep{Lovelock1971} given the Lean-verified LLC and causal invariance. The foliation-independence of the quantum measure
(causal invariance) is the RA statement of Lorentz covariance at the discrete
level, and its combinatorial core has now been formally verified in the Lean~4
theorem prover. Furthermore, by defining actualization via an irreversible
increase in quantum relative entropy and leveraging the Continuous Functional
Calculus, we prove the Stationarity Criterion for the Rindler thermal state,
resolving the Unruh paradox without breaking Lorentz covariance. The algorithm
makes a specific, near-term falsifiable prediction: a strict null result for
gravity-mediated entanglement experiments (the Bose-Marletto-Vedral protocol),
because the spacetime metric is updated strictly from actualized vertices and
cannot be placed in quantum superposition. Substantial progress on the
generative measure: the Benincasa-Dowker sign criterion
($S_{\mathrm{BDG}}(v) > 0$) provides the discrete actualization filter,
yielding a Poisson-CSG formula $c_k^{\mathrm{RA}}(x) \propto e^{-\mu}\mu^k$
with $\mu = \lambda V_{\mathrm{coh}} p_{\mathrm{th}}$ that is local and
environment-dependent, recovering $d_s = 3$ in dense environments and
$d_s = 2$ (Durhuus-Jonsson-Wheater) in sparse ones.

\medskip\noindent
\textbf{Keywords:} Quantum gravity \textperiodcentered{} Causal set theory
\textperiodcentered{} Quantum measure \textperiodcentered{} Causal invariance
\textperiodcentered{} Process ontology \textperiodcentered{}
Bose-Marletto-Vedral \textperiodcentered{} Relational Actualism
\textperiodcentered{} Benincasa-Dowker \textperiodcentered{} Poisson-CSG
\end{abstract}

\newpage

% ─────────────────────────────────────────────────────────────────────────────
\section{Introduction: The Shared Error}
\label{sec:intro}
% ─────────────────────────────────────────────────────────────────────────────

The reconciliation of quantum mechanics and general relativity has been the
central unsolved problem of theoretical physics for nearly a century. Every
proposed resolution---string theory, loop quantum gravity, causal dynamical
triangulations, asymptotic safety---begins by accepting both frameworks as given
and attempting to make them consistent. This paper proposes that the difficulty
is not in the reconciliation but in the starting point. Both frameworks share a
foundational error, and removing that error dissolves the tension rather than
resolving it.

The shared error is this: both quantum mechanics and general relativity treat
time as a reversible background parameter. In the Schr\"{o}dinger equation, $t$
is a continuous coordinate. In the Einstein field equations, $t$ is a coordinate
label on a block manifold. In both cases, \emph{nothing actually happens}. The
formalism describes a static mathematical structure; genuine temporal becoming
is represented nowhere in either framework.

Relational Actualism~\citep{Sandeman2026a} proposes that time is not a background
parameter but a physical process: the irreversible growth of the causal Directed
Acyclic Graph (DAG) of actualized events. Each vertex written into the DAG is an
on-shell actualization event---an irreversible increase in the entanglement entropy
of a system with its asymptotic environment, generating a redundant causal record.

This paper presents the \textbf{Engine of Becoming}: the complete formal algorithm
by which the causal graph grows. It is a Covariant Stochastic Graph Rewriting
System stated at three levels of description. The fundamental level is purely
graph-theoretic: no Hilbert space, no metric. The intermediate level introduces
the quantum measure---a foliation-independent assignment of complex amplitudes to
causal histories, due to Sorkin~\citep{Sorkin1994}. The emergent level recovers
the Feynman path integral and Einstein's field equations as continuum limits.

Quantum mechanics and general relativity are not rival fundamental theories. They
are descriptions of different levels and limits of the same algorithm.

% ─────────────────────────────────────────────────────────────────────────────
\section{Why Both Frameworks Fail}
\label{sec:failure}
% ─────────────────────────────────────────────────────────────────────────────

\subsection{The Schr\"{o}dinger Equation: Determinism Without Actuality}

The Schr\"{o}dinger equation is a deterministic, reversible, unitary evolution
law. It contains no mechanism by which a superposition of possibilities becomes
a definite actual outcome. The measurement problem is not about observers or
detectors; it is about the Schr\"{o}dinger equation's fundamental inability to
distinguish between mathematical potentia and physical actuality.

Every proposed solution (many worlds, pilot wave, spontaneous collapse) either
denies that collapse is real, adds a new postulate to produce it, or makes collapse
observer-relative. None addresses the root cause: the Schr\"{o}dinger equation is
the correct perturbative approximation to the quantum possibility space (Level~2
of the algorithm below), but it was never the fundamental dynamical law.

\subsection{The Einstein Field Equations: Geometry Without Process}

The Einstein field equations describe how the spacetime metric responds to
stress-energy. They are formulated on a differentiable manifold in which all of
spacetime, past and future, exists simultaneously. Quantizing this block is the
wrong approach for the same reason that quantizing fluid mechanics is wrong as
a path to atomic physics: the continuous manifold is itself the hydrodynamic
limit (Level~3 of the algorithm), not the fundamental object.

The tension between QM and GR is therefore a conflict between two different
approximations to the same underlying algorithm, applied in regimes where both
matter simultaneously.

% ─────────────────────────────────────────────────────────────────────────────
\section{The Covariance Problem and Its Resolution}
\label{sec:covariance}
% ─────────────────────────────────────────────────────────────────────────────

Before stating the algorithm, the central technical challenge must be addressed
directly. Any algorithm that selects actualization events in a preferred sequential
order---``step $n$, then step $n+1$''---is foliation-dependent: it implicitly
picks a preferred time slicing. Special relativity forbids this. Two
spacelike-separated actualization events have no objective temporal ordering.

This is the problem that prevented earlier collapse models (GRW and its
variants) from achieving Lorentz covariance, and it is the problem this paper
must resolve to deliver a genuine quantum gravity framework.

\subsection{Sorkin's Quantum Measure}

The resolution comes from Sorkin's quantum measure theory~\citep{Sorkin1994,
Sorkin2007}. A \emph{quantum measure} $\muC$ is a generalisation of a classical
probability measure defined on the space of complete histories (paths through
the causal DAG), such that:
\begin{enumerate}[noitemsep]
  \item Amplitudes are assigned to individual histories: $\mathcal{A}(h) \in \mathbb{C}$.
  \item The quantum measure of a set of histories $\mathcal{S}$ is the squared
  modulus of the sum of amplitudes:
  \begin{equation}
    \muC(\mathcal{S}) \;=\; \left|\sum_{h \in \mathcal{S}} \mathcal{A}(h)\right|^2.
    \label{eq:quantum_measure}
  \end{equation}
  \item Unlike classical probability, quantum measures allow interference
  between histories: $\muC(\mathcal{S}_1 \cup \mathcal{S}_2) \neq
  \muC(\mathcal{S}_1) + \muC(\mathcal{S}_2)$ in general.
\end{enumerate}
Classical probability is the special case where all amplitudes are
non-negative real numbers. Quantum mechanics is the case where amplitudes
are complex and interference is non-trivial.

\subsection{Causal Invariance: The Discrete Lorentz Covariance Condition}

The foliation-dependence problem is solved by the condition of \emph{causal
invariance}: the quantum measure must be independent of the order in which
spacelike-separated vertices are added to the causal DAG.

\begin{definition}[Causal Invariance]
A quantum measure $\muC$ on the growth process of a causal DAG $G$ is
\emph{causally invariant} if, for any two spacelike-separated candidate
vertices $v_A$ and $v_B$ (neither in the causal past of the other), and
for any partial graph $G_n$:
\begin{equation}
  \muC\bigl(G_n \to G_n \cup \{v_A, v_B\}\bigr)
  \;=\;
  \muC\bigl(G_n \to G_n \cup \{v_A\} \to G_n \cup \{v_A, v_B\}\bigr)
  \;=\;
  \muC\bigl(G_n \to G_n \cup \{v_B\} \to G_n \cup \{v_A, v_B\}\bigr).
  \label{eq:causal_invariance}
\end{equation}
\end{definition}

Causal invariance is the discrete analogue of Lorentz covariance: the physics
of spacelike-separated events does not depend on which is ``first.'' In the
continuum limit, causal invariance reduces to the microcausality condition of
quantum field theory: the commutator of spacelike-separated field operators
vanishes~\citep{Haag1996}.

Sorkin established causal invariance for the deterministic sequential growth
of causal sets~\citep{RideoutSorkin2000}. The primary hard wall of the Engine
of Becoming is the extension of this result to the quantum measure with
Born-rule weighting (Section~\ref{sec:hardwalls}).

% ─────────────────────────────────────────────────────────────────────────────
\section{The Engine of Becoming: Three Levels}
\label{sec:algorithm}
% ─────────────────────────────────────────────────────────────────────────────

The algorithm is stated at three levels of description. Level~1 is the
graph-theoretic primitive: no Hilbert space, no metric, only the causal DAG
and its growth rules. Level~2 introduces the quantum measure: complex amplitudes
on causal histories, foliation-independent by causal invariance. Level~3 is
the emergent limit: the Feynman path integral and Einstein equations appear
as continuum approximations to Levels~1 and~2 when the graph is sufficiently
dense and slowly varying.

\subsection{Level 1: The Graph-Theoretic Primitive}

Let the universe be defined by the causal DAG $G = (V, E)$, where $V$ is the
set of actualized vertices and $E$ the set of directed causal edges. Three
primitive concepts:

\begin{definition}[Causal Boundary]
The causal boundary $\partial G$ is the set of vertices in $V$ that have no
causal successors yet:
\begin{equation}
  \partial G \;=\; \{v \in V \mid \nexists\, w \in V : (v,w) \in E\}.
  \label{eq:boundary_v2}
\end{equation}
\end{definition}

\begin{definition}[Candidate Extension]
A candidate vertex $v_*$ is a potential new element of $V$ whose causal
antecedents are a non-empty subset of $\partial G$. The set of all candidate
extensions of $G$ is denoted $\mathcal{C}(G)$.
\end{definition}

\begin{definition}[Actualization Criterion]
A candidate $v_* \in \mathcal{C}(G)$ is \emph{valid} if it satisfies the
actualization criterion of RAQM~\citep{Sandeman2026a}: its actualization would
produce an irreversible increase in the quantum relative entropy of the system
with respect to the vacuum reference state, $\Delta S(\rho\|\sigma_0) > 0$.
In the perturbative regime, this coincides with the on-shell condition
$p^\mu p_\mu \geq m_b^2 c^2$ for the mediating field excitation. The valid
candidate set is $\Vvalid \subseteq \mathcal{C}(G)$.
\end{definition}

At Level~1, the universe grows by the addition of elements from $\Vvalid$ to
$V$. This is a purely combinatorial, background-independent process: no
spacetime, no Hilbert space. The past is the fixed set $V$; the future is the
set of all possible extensions $\mathcal{C}(G)$; the present is the boundary
$\partial G$.

\subsection{Level 2: The Quantum Measure}

Level~2 equips the Level~1 growth process with quantum amplitudes. A
\emph{causal history} $h$ is a maximal chain of single-element extensions
from the empty set to the complete DAG $G$. The quantum measure assigns:
\begin{equation}
  \mathcal{A}(h) \;=\; \prod_{v \in h} a(v \mid \mathrm{past}(v)),
  \label{eq:local_amplitude}
\end{equation}
where $a(v \mid \mathrm{past}(v))$ is the local amplitude for adding vertex
$v$ given its causal antecedents. The probability of any observable
(any set of partial DAGs $\mathcal{S}$) is:
\begin{equation}
  P(\mathcal{S}) \;=\; \muC(\mathcal{S})
  \;=\; \left|\sum_{h \supset \mathcal{S}} \mathcal{A}(h)\right|^2.
  \label{eq:qm_probability}
\end{equation}
Causal invariance (Definition~2) requires that $\muC(\mathcal{S})$ is
independent of the foliation of $h$ used to define the sequential ordering.

\paragraph{The local amplitude.}
In the perturbative regime, the local amplitude $a(v \mid \mathrm{past}(v))$
is the Feynman propagator~\citep{Feynman1948} between the causal antecedents of $v$ and $v$ itself.
This is not a separate postulate: the Feynman propagator is the unique
amplitude consistent with Lorentz covariance, unitarity, and the causal
structure of the DAG. The causal antecedents of $v$ define the boundary
conditions; the propagator computes the amplitude from those boundary conditions
to the candidate vertex. Quantum field theory is Level~2 in the limit where
the graph is dense enough to approximate a continuous spacetime background.

\paragraph{Wave function collapse without a preferred frame.}
In the covariant formulation, ``collapse'' is not the selection of one vertex
per step. It is the assignment of zero amplitude to all histories inconsistent
with the actualized subgraph. After any set of causally disconnected
actualizations $\{v_1, \ldots, v_k\}$ has been added to $V$, all candidate
histories that conflict with any $v_i$ are removed from the sum in
equation~\eqref{eq:qm_probability}. This is collapse without a preferred
ordering: the pruning is determined by the causal structure of $G$, not by
a clock.

\subsection{Level 3: Macroscopic Approximations}

Level~3 is what Levels~1 and~2 look like to observers who resolve
only large-scale averages. It comprises two regimes with distinct
epistemic status.

\paragraph{Quantum mechanics: exact at Level~2, approximated at Level~3.}
The quantum measure of Level~2 assigns complex amplitudes to causal histories
exactly, without approximation. Quantum mechanics is not a limit of the
Engine of Becoming; it IS the Engine of Becoming at Level~2. When the causal
graph is slowly varying, the quantum measure reduces to the Feynman path
integral:
\begin{equation}
  \muC(\text{transition: } |i\rangle \to |f\rangle)
  \;\xrightarrow{\mu \gg 1}\;
  \left|\int \mathcal{D}\Phi\; e^{iS[\Phi]/\hbar}\right|^2.
  \label{eq:qm_limit}
\end{equation}
The Schr\"{o}dinger equation emerges as the non-relativistic, single-particle
level of equation~\eqref{eq:qm_limit}. The Feynman integral is Level~3 QM;
Level~2 is more fundamental and exact.

\paragraph{General relativity: inevitable at Level~3, not fundamental.}
When the actualization density is high ($\mu \gg 1$) and the graph is
sufficiently isotropic, the Level~1 growth process produces a smooth
manifold whose geometry is Minkowski in the flat limit and Lorentzian in
the presence of matter. By Lovelock's theorem~\citep{Lovelock1971}, the
unique consistent covariant, conserving, local field equations on a
4-dimensional Lorentzian manifold are the Einstein equations. Since the
RA-CSG satisfies all three conditions (conservation: LLC, Lean-verified;
covariance: causal invariance, Lean-verified; locality: definitional in CSG),
the Einstein equations are not recovered in a limit --- they are logically
guaranteed. The interesting physics is where RA \emph{departs} from GR:
sparse regions ($\mu \sim 1$), Planck-scale corrections, Causal Severance.

The local actualization density:
\begin{equation}
  \lambda(x) \;=\; \lim_{\mathcal{V} \to \mathrm{meso}}
  \frac{|V \cap \mathcal{V}|}{\mathrm{Vol}(\mathcal{V})}
  \label{eq:lambda_density_v2}
\end{equation}
sources an effective metric via the RA-modified Einstein equations:
\begin{equation}
  G_{\mu\nu} \;=\; 8\pi G\,\langle\Tact\rangle,
  \label{eq:ra_efe}
\end{equation}
where $\langle\Tact\rangle$ is the RA-projected stress-energy restricted to
actualized on-shell states (RAGC~\citep{Sandeman2026b}). The Einstein equations are not fundamental; they are the
large-$\mu$ inevitable description of the Level~1 actualization
density, guaranteed by Lovelock. General relativity is the
hydrodynamic shadow of the causal graph at $\mu \gg 1$.

\paragraph{The relationship between levels.}
Level~1 determines what can actualize (the causal structure). Level~2
determines the probabilities (quantum interference). Level~3 is what Levels~1
and~2 look like to observers who can only resolve large-scale averages. The
Schr\"{o}dinger equation and the Einstein equations are both Level~3
descriptions. They are incompatible not because nature is inconsistent but
because they are approximations valid in different regimes of the same
underlying process, and they cannot both be fundamental.

% ─────────────────────────────────────────────────────────────────────────────
\section{The Five Steps Revisited}
\label{sec:fivesteps}
% ─────────────────────────────────────────────────────────────────────────────

The five-step description of the algorithm can now be restated in manifestly
covariant form, with each step identified with its level.

\begin{description}[leftmargin=1.5em, itemsep=6pt]
  \item[\textbf{Step 1 --- The Causal Boundary (Level~1).}]
  The present is $\partial G$: the set of actualized vertices with no
  causal successors. This is a causal notion, not a spatial one.
  Different observers have different causal pasts and therefore different
  $\partial G$, but the Level~1 growth rules are the same for all observers.

  \item[\textbf{Step 2 --- The Quantum Possibility Space (Levels~1--2).}]
  From $\partial G$, the valid candidate set $\Vvalid$ is determined by the
  actualization criterion (Level~1). Amplitudes are assigned to all possible
  extensions via the quantum measure (Level~2). Virtual candidates carry zero
  metric weight (they are not in $V$) and do not source gravity. This is the
  structural explanation for why the vacuum does not gravitate.

  \item[\textbf{Step 3 --- The Kinematic Snap (Level~1).}]
  The actualization criterion filters $\mathcal{C}(G)$ to $\Vvalid$ by
  requiring an irreversible increase in quantum relative entropy
  $\Delta S(\rho\|\sigma_0) > 0$. Because Lorentz boosts are represented
  by unitary conjugations, and the Continuous Functional Calculus guarantees
  that relative entropy is unitarily invariant~\citep{Sandeman2026a}, this
  threshold is an objective Lorentz scalar. A uniformly accelerated (Rindler)
  observer sees a KMS thermal bath, but its relative entropy with respect to
  the vacuum is constant under modular flow ($\Delta S_{\mathrm{rel}} = 0$,
  proved in \texttt{RA\_AQFT\_Proofs\_v10.lean}): no actualization occurs.
  The Unruh effect is thereby resolved without breaking Lorentz covariance---
  the Rindler observer agrees with the inertial observer that no physical snap
  has occurred, because neither observer's $\Delta S$ threshold is crossed.

  \item[\textbf{Step 4 --- The Covariant Stochastic Update (Level~2).}]
  A set of causally disconnected valid candidates $\{v_1, \ldots, v_k\} \in
  \Vvalid$ is actualized with probability given by the quantum measure:
  \begin{equation}
    P(\{v_1,\ldots,v_k\})
    \;=\; \left|\sum_\sigma \prod_i a(v_{\sigma(i)}\mid
      \mathrm{past}(v_{\sigma(i)}))\right|^2,
    \label{eq:covariant_update}
  \end{equation}
  where the sum is over all orderings $\sigma$ of $\{v_1,\ldots,v_k\}$
  consistent with the causal partial order. By causal invariance, this is
  the same for every foliation: the probability depends only on the causal
  structure of the antecedents, not on the sequential labelling. This is
  the covariant replacement for the foliation-dependent ``select one vertex
  per step.''

  \item[\textbf{Step 5 --- The Hydrodynamic Limit (Level~3).}]
  After actualization, $V$ is updated and the effective metric is
  recalculated from the new actualization density $\lambda(x)$ via
  equation~\eqref{eq:ra_efe}. This is a Step~3 output: it happens only
  after actualization, never during the quantum exploration. Gravity is
  a post-actualization recalculation. Superposed states never source the
  metric.
\end{description}

% ─────────────────────────────────────────────────────────────────────────────
\section{The BMV Prediction}
\label{sec:bmv}
% ─────────────────────────────────────────────────────────────────────────────

\subsection{The Bose-Marletto-Vedral Protocol}

Bose et al.~\citep{Bose2017} and Marletto and Vedral~\citep{Marletto2017}
proposed that quantum effects in gravity could be detected by observing
gravity-mediated entanglement between two massive particles held in spatial
superpositions. If gravity mediates entanglement between superposed systems,
gravity must be quantum-mechanical. Multiple experimental groups are actively
pursuing this test.

\subsection{The RA Prediction: Strict Null Result}

The Engine of Becoming predicts a strict null result: no gravity-mediated
entanglement between superposed masses will be observed.

The argument follows directly from Step~5. The spacetime metric is recalculated
strictly from $V$---the set of actualized vertices---and never from the quantum
possibility space (Step~2). A superposed mass that has not yet actualized
contributes zero weight to $V$ and therefore zero metric sourcing. The metric
cannot be placed in superposition because the metric is a property of the
actualized causal record, not of the space of unactualized possibilities.

\paragraph{Distinction from semiclassical gravity.}
Standard semiclassical gravity sources the metric from the expectation value
$\langle T_{\mu\nu}\rangle$, which is non-zero for a superposed mass. RA
sources it from $\langle\Tact\rangle$, which is zero for an unactualized
superposition. The predictions are quantitatively distinct and experimentally
discriminable.

\paragraph{What a positive BMV result would mean.}
A positive result would falsify the Step~2/Step~5 separation: gravity would
have to be sensitive to unactualized virtual states. This would require
fundamentally modifying the architecture in which the metric is updated only
from actualized vertices---a falsification of the Engine of Becoming as stated.

% ─────────────────────────────────────────────────────────────────────────────
\section{Recovery of QM and GR as Limiting Cases}
\label{sec:limits}
% ─────────────────────────────────────────────────────────────────────────────

\subsection{Quantum Mechanics as the Level~2 / Slowly-Varying Limit}

When the metric changes negligibly between actualization events, the quantum
measure reduces to the Feynman path integral on a fixed spacetime background.
This is quantum field theory. The Schr\"{o}dinger equation is the
non-relativistic, single-particle limit. Quantum mechanics is therefore the
correct description of Level~2 in the approximation that Level~3 changes slowly.

The measurement problem disappears because Level~2 provides the space of
possibilities and Level~4 (Step~4) provides the objective, covariant
actualization mechanism. The problem only appears when the Schr\"{o}dinger
equation is treated as the complete fundamental law, omitting Steps~3--4.

\subsection{General Relativity as the Level~3 / Dense-Graph Limit}

When the actualization rate is high and the graph is isotropic, Level~1 growth
produces a smooth manifold and the effective metric satisfies the Einstein
equations. GR is the correct description when the discrete actualization
structure averages into a smooth continuum. In the regimes where $\lambda$ is
inhomogeneous---galactic boundaries, voids, cluster collisions---the RA
corrections become significant and produce the testable deviations of
RAGC~\citep{Sandeman2026b} and RADM~\citep{Sandeman2026RADM}.

\subsection{The Unruh Paradox Resolved}

The Unruh effect is often cited as evidence that particle content is
observer-dependent, threatening the objectivity of actualization. In the RA
framework, the resolution is direct and Lean-verified.

A uniformly accelerated (Rindler) observer perceives the Minkowski vacuum as
a thermal KMS state at temperature $T_U = \hbar a / 2\pi c k_B$. Naively,
this suggests the Rindler observer should experience actualization events not
seen by the inertial observer---a violation of the objectivity of the
Kinematic Snap.

The resolution: the actualization criterion is $\Delta S(\rho\|\sigma_0) > 0$,
not particle number. The Rindler thermal state has $\Delta S_{\mathrm{rel}} = 0$
with respect to the Minkowski vacuum under modular flow---the relative entropy
is stationary. This is the content of the \texttt{rindler\_stationarity}
theorem proved in \texttt{RA\_AQFT\_Proofs\_v10.lean}: the relative entropy
between the Rindler KMS state and the vacuum is invariant under the modular
automorphism group of the Rindler wedge algebra. Because the threshold
$\Delta S > 0$ is not crossed, no actualization is attributed to the thermal
bath by the Rindler observer. Both observers agree: no physical snap has
occurred. The apparent paradox dissolves because the actualization criterion
is covariant by construction---relative entropy is unitarily invariant under
Lorentz boosts (proved via the Continuous Functional Calculus in
\texttt{frame\_independence}).

% ─────────────────────────────────────────────────────────────────────────────
\section{The Algorithm and the Suite}
\label{sec:suite}
% ─────────────────────────────────────────────────────────────────────────────

Each companion paper in the RA suite derives consequences of one or more
levels or steps of the Engine of Becoming:

\begin{description}[leftmargin=1.5em]
  \item[\textbf{RAQM}~\citep{Sandeman2026a}] Establishes the actualization
  criterion (Step~3), derives the Born rule (quantum measure on the graph),
  and solves the measurement problem.
  \item[\textbf{RAGC}~\citep{Sandeman2026b}] Derives the Local Ledger
  Condition (Step~4 pruning), vacuum energy suppression (Step~2 non-sourcing),
  and GR as the Level~3 hydrodynamic limit (Step~5).
  \item[\textbf{RASM}~\citep{Sandeman2026RASM}] Derives the Standard Model
  from BDG topology: charge quantisation, three generations, force structure,
  Koide formula, and the Wyler formula from the same causal diamond geometry.
  \item[\textbf{RADM}~\citep{Sandeman2026RADM}] Derives dark matter predictions
  from Step~5 in the inhomogeneous-$\lambda$ regime.
  \item[\textbf{RAQI}~\citep{Sandeman2026d}] Derives the Kinematic Coherence
  Bound from the Step~3 threshold.
  \item[\textbf{RAHC}~\citep{Sandeman2026RAHC}] Derives the full complexity
  hierarchy as repeated application of Step~5 at successively larger scales.
  \item[\textbf{RACI}~\citep{Sandeman2026c}] Applies Steps~3--5 to biological
  complexity, life, and intelligence.
\end{description}

% ─────────────────────────────────────────────────────────────────────────────
\section{Hard Walls and Open Problems}
\label{sec:hardwalls}
% ─────────────────────────────────────────────────────────────────────────────

\begin{description}
  \item[Hard Wall 1: Causal Invariance --- Status: Essentially Proved in the
  Perturbative Regime; Discrete Intrinsic Formulation Remains Open.]

  We present the proof of causal invariance in full for the perturbative
  regime, reducing the remaining open problem to a single precisely stated
  discrete locality condition.

  \paragraph{The Locality Lemma.}
  The entire proof rests on one lemma:

  \begin{definition}[Amplitude Locality]
  The local amplitude $a(v \mid S)$ is \emph{local} if it depends only on
  the elements of $S$ that lie in the causal past of $v$:
  \begin{equation}
    a(v \mid S) = a(v \mid S') \quad\text{whenever}\quad
    S \cap \mathrm{past}(v) = S' \cap \mathrm{past}(v).
    \label{eq:locality}
  \end{equation}
  \end{definition}

  Amplitude locality is the quantum analogue of Rideout and Sorkin's Bell
  causality condition~\citep{RideoutSorkin2000}, which requires that the
  classical transition \emph{probability} depends only on the causal past of
  the new element. Our condition lifts this from probabilities to complex
  amplitudes. In the continuum/perturbative limit, amplitude locality holds
  because the Feynman propagator between causal antecedents and $v$ is a
  Green's function of the wave equation, and signals propagate at most at $c$:
  a spacelike-separated element of $S$ cannot affect the boundary conditions
  relevant to $v$'s amplitude. This is the content of the Huygens-Fresnel
  principle and the vanishing of the Pauli-Jordan function for spacelike
  separations. In the operator language of QFT: the local amplitude is a
  matrix element of $\hat{\phi}(v)$; the microcausality theorem
  $[\hat{\phi}(x), \hat{\phi}(y)] = 0$ for spacelike $(x,y)$~\citep{Haag1996}
  guarantees that spacelike elements of $S$ commute through the matrix element
  and leave $a(v \mid S)$ unchanged.

  \paragraph{Theorem (Causal Invariance).}
  \begin{theorem}[RA Causal Invariance]
  \label{thm:causal_invariance}
  Let $\{v_1, \ldots, v_N\}$ be a set of mutually spacelike valid candidates.
  If the local amplitudes are amplitude-local in the sense of
  equation~\eqref{eq:locality}, then the quantum measure
  $\muC(\{v_1,\ldots,v_N\})$ is independent of the ordering $\sigma$:
  \begin{equation}
    \muC = \bigl|\mathcal{A}(\sigma)\bigr|^2
    \quad\text{is the same for all orderings } \sigma.
    \label{eq:ci_theorem}
  \end{equation}
  \end{theorem}

  \begin{proof}
  \textbf{Base case ($N=2$).} Let $v_A, v_B$ be spacelike, with causal pasts
  $P_A = \mathrm{past}(v_A)$ and $P_B = \mathrm{past}(v_B)$ in $G_n$. Since
  $v_A \notin P_B$ and $v_B \notin P_A$ (spacelike separation), amplitude
  locality gives:
  \begin{align}
    a(v_B \mid P_B \cup \{v_A\}) &= a(v_B \mid P_B), \label{eq:loc_B}\\
    a(v_A \mid P_A \cup \{v_B\}) &= a(v_A \mid P_A). \label{eq:loc_A}
  \end{align}
  The history amplitudes for the two orderings are therefore:
  \begin{align}
    \mathcal{A}(v_A \text{ then } v_B)
      &= a(v_A \mid P_A) \cdot a(v_B \mid P_B \cup \{v_A\})
       = a(v_A \mid P_A) \cdot a(v_B \mid P_B), \\
    \mathcal{A}(v_B \text{ then } v_A)
      &= a(v_B \mid P_B) \cdot a(v_A \mid P_A \cup \{v_B\})
       = a(v_B \mid P_B) \cdot a(v_A \mid P_A).
  \end{align}
  Since multiplication of complex numbers is commutative:
  $\mathcal{A}(v_A \text{ then } v_B) = \mathcal{A}(v_B \text{ then } v_A)$.
  Therefore $\muC = |\mathcal{A}|^2$ is identical for both orderings.

  \textbf{Inductive step.} Assume the result holds for all mutually spacelike
  sets of size $< N$. Every permutation of $N$ elements is a composition of
  adjacent transpositions (swaps of elements at positions $i$ and $i+1$). It
  therefore suffices to show that swapping adjacent elements $v_i$ and $v_{i+1}$
  (both spacelike to each other) leaves $\mathcal{A}(\sigma)$ unchanged.

  By amplitude locality, the prefix $\prod_{j<i} a(v_{\sigma(j)} \mid \cdots)$
  and suffix $\prod_{j>i+1} a(v_{\sigma(j)} \mid \cdots)$ are unaffected by the
  swap (they involve vertices whose causal pasts do not include $v_i$ or
  $v_{i+1}$ by spacelike separation). For the swapped pair, the same two-element
  argument applies:
  \begin{equation}
    a(v_i \mid P_i) \cdot a(v_{i+1} \mid P_{i+1})
    = a(v_{i+1} \mid P_{i+1}) \cdot a(v_i \mid P_i),
  \end{equation}
  by commutativity of $\mathbb{C}$. Therefore $\mathcal{A}(\sigma_1) =
  \mathcal{A}(\sigma_2)$ for any two orderings $\sigma_1, \sigma_2$, and
  $\muC = |\mathcal{A}(\sigma)|^2$ is ordering-independent. \qed
  \end{proof}

  \paragraph{Status and remaining gap.}
  The proof is complete in the perturbative/continuum regime where the local
  amplitude $a(v \mid \mathrm{past}(v))$ is identified with the Feynman
  propagator. In that regime, amplitude locality is a theorem, not an
  assumption---it follows from QFT microcausality~\citep{Haag1996} and the
  Pauli-Jordan function vanishing outside the light cone. The main theorem
  then follows by elementary algebra.

  The remaining open problem---now precisely stated---is:
  \begin{quotation}
  \noindent\textit{Define the local amplitude $a(v \mid \mathrm{past}(v))$
  intrinsically on the discrete RA-CSG graph, without appealing to the
  continuum limit, and prove that it satisfies amplitude locality
  (equation~\eqref{eq:locality}) as a theorem of the discrete graph dynamics
  rather than as an inherited property of the QFT approximation.}
  \end{quotation}
  \textbf{Lean 4 verification status: complete (given amplitude locality).}
  The combinatorial structure of this theorem---the induction on adjacent
  transpositions, the foliation-independence of \texttt{Finset.prod} over a
  \texttt{CommMonoid}, and the causal partial-order properties---has been
  formally verified in \texttt{RA\_AQFT\_Proofs\_v2.lean} with zero messages:
  no errors, no warnings, no unresolved goals. The file is available at
  \href{https://github.com/jsandeman/Relational-Actualism}{github.com/jsandeman/Relational-Actualism}.
  The sole remaining open problem is the intrinsic discrete proof of the
  \texttt{amplitude\_locality} axiom stated above---showing that the local
  amplitude on the discrete RA-CSG graph satisfies locality as a theorem
  of the discrete dynamics rather than as an inherited property of the QFT
  continuum limit.
  \textbf{Collaborator target:} causal set theory~\citep{Dowker2005,RideoutSorkin2000} (Dowker, Sorkin, Henson).

  \item[Hard Wall 2: Derivation of Local Amplitudes and the Generative Measure.]
  The local amplitude $a(v \mid \mathrm{past}(v))$ is currently taken from
  the Feynman propagator---a Level~3 object used at Level~2. A complete
  derivation must show that the propagator \emph{emerges} from the graph
  structure of $G$ without being postulated.

  Substantial progress has been made on the closely related measure problem.
  The Benincasa-Dowker-Glaser (BDG) action~\citep{Benincasa2010} provides the
  unique discrete, retarded, Lorentz-invariant action functional for a causal
  set (pure graph combinatorics, no metrics). The RA-native actualization filter
  $\Delta S(\rho\|\sigma_0) > 0$ has a natural discrete counterpart: a candidate
  vertex $v$ is actualized if and only if $S_{\mathrm{BDG}}(v) > 0$, identifying
  a positive contribution to local scalar curvature with a positive increase in
  relative entropy. This BDG sign criterion, combined with QFT cluster
  decomposition (spacelike-separated interactions decohere independently), yields
  the Poisson-CSG generative measure~\citep{Sandeman2026b}:
  \begin{equation}
    c_k^{\mathrm{RA}}(x) = \frac{e^{-\mu(x)}\,\mu(x)^k}{k!},
    \qquad
    \mu(x) = \lambda(x)\,V_{\mathrm{coh}}(x)\,p_{\mathrm{th}},
  \end{equation}
  where $\lambda(x)$ is the local actualization density, $V_{\mathrm{coh}}$ is
  the pre-actualization coherence volume, and $p_{\mathrm{th}} \approx 10^{-2}$
  is the Kinematic Coherence Bound threshold of RAQI~\citep{Sandeman2026d}.
  This makes $c_k$ \emph{local and environment-dependent} for the first time in
  any CSG formulation. In dense environments ($\mu \gg 1$): $d_s = 3$, GR is
  recovered. In sparse environments ($\mu \ll 1$): the graph becomes tree-like
  and the Durhuus-Jonsson-Wheater theorem~\citep{Durhuus2010} gives $d_s = 2$,
  driving dimensional reduction and the conditional Tully-Fisher derivation of
  RADM~\citep{Sandeman2026RADM}. The remaining gap is determining $V_{\mathrm{coh}}$
  from the vacuum state structure and running the BDG-filter MCMC to extract
  the full asymptotic $c_k$ distribution.
  \textbf{Collaborator target:} causal set theory, CDT computational physics,
  algebraic QFT~\citep{Dowker2005,RideoutSorkin2000}.

  \item[Hard Wall 3: The Standard Model Couplings.]
  The local amplitudes depend on the coupling constants of the Standard Model.
  The Engine of Becoming uses these as inputs rather than deriving them from
  the graph topology. The derivation of $U(1) \times SU(2) \times SU(3)$ from
  the causal DAG structure is the long-range horizon described in the RAGC
  Epilogue~\citep{Sandeman2026b}.
  \textbf{Collaborator target:} topological quantum field theory, Standard
  Model phenomenology.
\end{description}

% ─────────────────────────────────────────────────────────────────────────────
\section{Conclusion}
\label{sec:conclusion}
% ─────────────────────────────────────────────────────────────────────────────

The Engine of Becoming is the formal answer to what the universe is doing,
step by step, as the causal graph grows. At Level~1, the graph-theoretic
primitive: candidates are identified by causal structure alone. At Level~2,
the quantum measure: complex amplitudes are assigned to causal histories in a
foliation-independent way. At Level~3, the emergent limit: the Feynman path integral appears as the large-$\mu$ macroscopic
approximation to the quantum measure; the Einstein field equations are
the unique consistent Level~3 description, inevitable by Lovelock's
theorem given the Lean-verified conservation and covariance.

Quantum mechanics and general relativity are not rival theories of nature.
They are descriptions of different levels of the same algorithm. The
measurement problem disappears because collapse is the covariant actualization
of one set of causally consistent candidates, not the mysterious selection of
one branch of a universal wave function. The quantum gravity problem
disappears because there is nothing to quantize: the metric is a
post-actualization coarse-graining of the causal graph, not a quantum field.

The combinatorial core of causal invariance has been formally verified in
Lean~4 (\texttt{RA\_AQFT\_Proofs\_v2.lean}, zero messages). The central
remaining open problem is Hard Wall~1 in its precise form: the intrinsic
discrete proof of amplitude locality, showing that the local amplitude on the
RA-CSG graph satisfies causal independence as a theorem of the discrete
dynamics, not merely as an inherited property of the QFT continuum limit.
This is the remaining step that would constitute the formal completion of the
Engine of Becoming.

The BMV experiment is the near-term experimental test. The Engine of Becoming
predicts a strict null result. If gravity-mediated entanglement between
superposed masses is observed, the Step~2/Step~5 separation is falsified and
the framework must be fundamentally revised. If no such entanglement is
observed, the experiment provides the first direct evidence that spacetime is
not a quantum field to be quantized but a hydrodynamic average over an
underlying discrete causal process.

The universe is a relentlessly self-writing causal history.
The Engine of Becoming is its algorithm.

% ─────────────────────────────────────────────────────────────────────────────

\section*{Declarations}
\textbf{Funding:} Not applicable.\quad
\textbf{Competing interests:} The author declares no competing interests.

\textbf{Preprint availability:} This manuscript is available as a preprint at
\href{https://doi.org/10.5281/zenodo.19198120}{doi.org/10.5281/zenodo.19198120}.

\textbf{Code availability:} The Lean~4 formal proof files
(\texttt{RA\_AQFT\_Proofs\_v10.lean} and \texttt{RA\_AQFT\_Proofs\_v2.lean})
are publicly available at
\href{https://github.com/jsandeman/Relational-Actualism}{github.com/jsandeman/Relational-Actualism}.

\textbf{Data availability:} This paper presents theoretical work. No datasets
were generated or analysed.


% ─────────────────────────────────────────────────────────────────────────────
\bibliographystyle{plainnat}
\begin{thebibliography}{99}

\bibitem[Sandeman(2026)]{Sandeman2026a}
J.~F.~Sandeman.
\newblock Relational Actualism and Quantum Mechanics ({RAQM}).
\newblock \textit{Foundations of Physics} (Under review), 2026.
\newblock \href{https://doi.org/10.5281/zenodo.19174380}{doi.org/10.5281/zenodo.19174380}.

\bibitem[Sandeman(2026)]{Sandeman2026b}
J.~F.~Sandeman.
\newblock Relational Actualism: Gravity and Cosmology ({RAGC}).
\newblock Preprint, 2026.
\newblock \href{https://doi.org/10.5281/zenodo.19197999}{doi.org/10.5281/zenodo.19197999}.

\bibitem[Sandeman(2026)]{Sandeman2026c}
J.~F.~Sandeman.
\newblock Relational Actualism: Complexity and Intelligence ({RACI}).
\newblock Preprint, 2026.
\newblock \href{https://doi.org/10.5281/zenodo.19198224}{doi.org/10.5281/zenodo.19198224}.

\bibitem[Sandeman(2026)]{Sandeman2026d}
J.~F.~Sandeman.
\newblock Relational Actualism: Implications for Quantum Information ({RAQI}).
\newblock Preprint, 2026.
\newblock \href{https://doi.org/10.5281/zenodo.19198285}{doi.org/10.5281/zenodo.19198285}.

\bibitem[Sandeman(2026)]{Sandeman2026RAHC}
J.~F.~Sandeman.
\newblock Algorithmic Causal Dynamics: A Unified Framework for Emergence
  and Complexity in Relational Actualism ({RAHC}).
\newblock Preprint, 2026.
\newblock \href{https://doi.org/10.5281/zenodo.19198171}{doi.org/10.5281/zenodo.19198171}.

\bibitem[Sandeman(2026)]{Sandeman2026RADM}
J.~F.~Sandeman.
\newblock Relational Actualism: Dark Matter as Actualization-Density
  Topology ({RADM}).
\newblock Preprint, 2026.
\newblock \href{https://doi.org/10.5281/zenodo.19198513}{doi.org/10.5281/zenodo.19198513}.

\bibitem[Lovelock(1971)]{Lovelock1971}
D.~Lovelock.
\newblock The Einstein tensor and its generalizations.
\newblock \textit{Journal of Mathematical Physics}, 12:498--501, 1971.

\bibitem[Sorkin(1994)]{Sorkin1994}
R.~D.~Sorkin.
\newblock Quantum mechanics as quantum measure theory.
\newblock \textit{Modern Physics Letters A}, 9(33):3119--3127, 1994.

\bibitem[Sorkin(2007)]{Sorkin2007}
R.~D.~Sorkin.
\newblock An exercise in ``anhomomorphic logic''.
\newblock \textit{Journal of Physics: Conference Series}, 67:012018, 2007.

\bibitem[Rideout and Sorkin(2000)]{RideoutSorkin2000}
D.~P.~Rideout and R.~D.~Sorkin.
\newblock Classical sequential growth dynamics for causal sets.
\newblock \textit{Physical Review D}, 61:024002, 2000.

\bibitem[Haag(1996)]{Haag1996}
R.~Haag.
\newblock \textit{Local Quantum Physics: Fields, Particles, Algebras}.
\newblock Springer, 2nd edition, 1996.

\bibitem[Bose et~al.(2017)]{Bose2017}
S.~Bose, A.~Mazumdar, G.~W.~Morley, H.~Ulbricht, M.~Toro\v{s},
  M.~Paternostro, A.~A.~Geraci, P.~F.~Barker, M.~S.~Kim, and G.~Milburn.
\newblock Spin entanglement witness for quantum gravity.
\newblock \textit{Physical Review Letters}, 119:240401, 2017.

\bibitem[Marletto and Vedral(2017)]{Marletto2017}
C.~Marletto and V.~Vedral.
\newblock Gravitationally induced entanglement between two massive particles is
  sufficient evidence of quantum effects in gravity.
\newblock \textit{Physical Review Letters}, 119:240402, 2017.


\bibitem[Benincasa \& Dowker(2010)]{Benincasa2010}
D.~M.~T.~Benincasa and F.~Dowker.
\newblock The scalar curvature of a causal set.
\newblock \emph{Physical Review Letters}, 104(18):181301, 2010.
\newblock \href{https://doi.org/10.1103/PhysRevLett.104.181301}{doi:10.1103/PhysRevLett.104.181301}.

\bibitem[Durhuus et~al.(2010)]{Durhuus2010}
B.~Durhuus, T.~Jonsson, and J.~F.~Wheater.
\newblock On the spectral dimension of causal triangulations.
\newblock \emph{Journal of Statistical Physics}, 139:859--881, 2010.
\newblock \href{https://doi.org/10.1007/s10955-010-9968-x}{doi:10.1007/s10955-010-9968-x}.

\bibitem[Dowker(2005)]{Dowker2005}
F.~Dowker.
\newblock Causal sets and the deep structure of spacetime.
\newblock In A.~Ashtekar, editor, \textit{100 Years of Relativity}, pages
  445--464. World Scientific, 2005.

\bibitem[Feynman(1948)]{Feynman1948}
R.~P.~Feynman.
\newblock Space-time approach to non-relativistic quantum mechanics.
\newblock \textit{Reviews of Modern Physics}, 20:367--387, 1948.

\bibitem[Sandeman(2026)]{Sandeman2026RASM}
J.~F.~Sandeman.
\newblock Relational Actualism and the Standard Model ({RASM}).
\newblock Preprint, 2026.
\newblock \href{https://doi.org/10.5281/zenodo.19229335}{doi.org/10.5281/zenodo.19229335}.

\bibitem[Sandeman(2026)]{Sandeman2026Foundation}
J.~F.~Sandeman.
\newblock Relational Actualism: A Foundation Paper.
\newblock Preprint, 2026.
\newblock \href{https://doi.org/10.5281/zenodo.19229438}{doi.org/10.5281/zenodo.19229438}.

\end{thebibliography}

\end{document}
